\documentclass[11pt,letterpaper]{article}
\usepackage[margin=0.60in,top=0.58in,bottom=0.56in]{geometry}
\usepackage{mathptmx}
\usepackage{microtype}
\usepackage{graphicx}
\usepackage{booktabs}
\usepackage{array}
\usepackage{longtable}
\usepackage{amsmath,amssymb}
\usepackage[hidelinks]{hyperref}
\setlength{\parindent}{0pt}
\setlength{\parskip}{2.5pt}
\setlength{\columnsep}{0.25in}
\setcounter{secnumdepth}{0}
\renewcommand{\section}[1]{\vspace{3pt}\noindent\textbf{\underline{#1}}\par\vspace{1pt}}
\renewcommand{\subsection}[1]{\vspace{2pt}\noindent\textbf{#1}\par\vspace{1pt}}
\renewcommand{\familydefault}{\rmdefault}
\providecommand{\tightlist}{\setlength{\itemsep}{0pt}\setlength{\parskip}{0pt}}
\providecommand{\pandocbounded}[1]{\sbox0{#1}\ifdim\wd0>\linewidth\resizebox{\linewidth}{!}{#1}\else\usebox0\fi}
\title{\vspace{-1.2em}\textbf{\MakeUppercase{Neurologically Motivated
Self-Supervised Learning for Bilateral Gait Geometry}}\vspace{-0.5em}}
\author{\normalsize [Authors and affiliations to be inserted in the RISEx template]}
\date{}
\begin{document}
\twocolumn[
  \begin{@twocolumnfalse}
  \maketitle
  \vspace{-1.1em}
  \end{@twocolumnfalse}
]
\subsection{INTRODUCTION}\label{introduction}

Robots that work around people need compact descriptions of human
motion. A useful description should retain geometric relations that
matter to a later task, rather than only make hidden inputs easy to
predict. Human gait offers one such relation. Reflecting a pose
exchanges left and right joints and reverses a signed difference between
their movement speeds. Gait asymmetry is clinically relevant in stroke
and Parkinson's disease, although the score studied here is not a
diagnosis {[}1{]}. We ask whether a skeleton JEPA retains this bilateral
movement signal after self-supervised training.

\subsection{MATERIALS AND METHODS}\label{materials-and-methods}

We used 625 accepted pose clips from 93 GAVD source videos {[}2{]}. Each
clip contains 33 landmarks. After within-clip normalization and
resizing, we retain all landmarks in a 64-by-33-by-3 input; four-frame
patches form a 16-by-33 grid. A JEPA encoder receives visible tokens and
predicts features of masked tokens supplied by a slowly updated teacher
encoder {[}3{]}. The signed target averages normalized left-minus-right
median speeds for shoulders, knees, ankles, heels, and foot tips.

The null hypothesis was that training would not improve a frozen
ridge-regression readout over the same encoder at initialization. We use
five outer source-video folds and five training seeds. Test videos are
excluded from encoder and readout fitting. Matched arms keep initial
weights, source draws, views, and update counts fixed while changing the
mask.

\subsection{RESULTS AND DISCUSSION}\label{results-and-discussion}

{\def\LTcaptype{none} % do not increment counter
\begin{tabular}{@{}p{0.43\columnwidth}p{0.47\columnwidth}@{}}
\toprule\noalign{}
\begin{minipage}[b]{\linewidth}\raggedright
Test
\end{minipage} & \begin{minipage}[b]{\linewidth}\raggedright
Result
\end{minipage} \\
\midrule\noalign{}
Initial frozen encoder & \(R^2=0.223\) \\
Trained teacher encoders & \(R^2=0.101\)--\(0.114\) \\
Trained minus initial & \(-0.109\) to \(-0.122\), five arms \\
Correct-clip feature target preferred & 375/375 trained checks; 33/75
initial checks \\
\end{tabular}
}

Training improved a diagnostic of within-clip feature correspondence
while reducing accuracy for the bilateral readout. Motion-weighted and
connected-region masks did not resolve this: their differences from
count-matched random masks had intervals that crossed zero. The result
is relevant to AI systems that represent articulated people, because
success on a latent prediction task did not imply preservation of this
downstream geometric quantity.

\subsection{CONCLUSIONS}\label{conclusions}

For the tested data, JEPA recipe, and frozen readout, hidden-feature
prediction and bilateral movement recovery disagree. The study provides
a compact diagnostic for evaluating learned articulated-motion
representations. It does not evaluate a robot, forecast future motion,
or establish clinical utility. Future work should test a past-only
human-motion prediction task with simple motion baselines and
independently specified data.

\subsection{REFERENCES}\label{references}

{[}1{]} Patterson KK et al.~\emph{Arch Phys Med Rehabil} 89:304--310,
2008. {[}2{]} Ranjan R et al.~\emph{IEEE Access} 13:45321--45339, 2025.
{[}3{]} Assran M et al.~arXiv:2301.08243, 2023.
\end{document}
